\documentclass[11pt,a4paper]{article}
% Formulação u - p_w - p_a com Mohr-Coulomb associado (Fase 1). Base: notas biot_nao_saturado_agua_ar3.tex e
% artigo Lira Cecílio (2026, J. Eng. Math. 157:10); cada equação confrontada com o código do NeoPZ.
% Compilar: pdflatex formulacao && pdflatex formulacao
\usepackage[T1]{fontenc}
\usepackage[utf8]{inputenc}
\usepackage[brazil]{babel}
\usepackage{lmodern}
\usepackage{amsmath,amssymb,bm,mathtools}
\usepackage{geometry}
\usepackage{booktabs,array,longtable,multirow}
\usepackage{xcolor}
\usepackage[hidelinks]{hyperref}
\geometry{margin=2.2cm}
\setlength{\parindent}{0pt}
\setlength{\parskip}{4pt}
\allowdisplaybreaks
\numberwithin{equation}{section}
\sloppy\setlength{\emergencystretch}{2em}

\newcommand{\Div}{\operatorname{div}}
\newcommand{\dd}{\mathrm{d}}
\newcommand{\Id}{\mathbf I}
\newcommand{\eps}{\boldsymbol\varepsilon}
\newcommand{\sig}{\boldsymbol\sigma}
\newcommand{\pabs}{p_a^{\mathrm{abs}}}
\newcommand{\sn}{\sin\phi}
\newcommand{\cs}{\cos\phi}
\newcommand{\code}[1]{\texttt{#1}}
% fontes: notas [N] e artigo [A]
\newcommand{\nota}[1]{{\footnotesize\textsf{[N~#1]}}}
\newcommand{\art}[1]{{\footnotesize\textsf{[A~#1]}}}
% trecho do código: arquivo e linhas (branch claude/relaxed-volta-1r4gog, 30/10/2026)
\newcommand{\loc}[2]{{\footnotesize[\path{#1}:#2]}}
% discrepância entre a formulação e o código (lista na Seção \ref{sec:disc})
\newcommand{\disc}[1]{\textcolor{red!70!black}{\textbf{\footnotesize[#1]}}}

\title{Formulação acoplada $\mathbf u$--$p_w$--$p_a$ com Mohr--Coulomb associado\\[3pt]
\large equações, discretização, confronto com o que o NeoPZ resolve e plano de verificação}
\author{Documento de trabalho, Fase 1}
\date{30 de outubro de 2026}

\begin{document}
\maketitle
\vspace{-1.5em}

\section*{Resumo}

Este documento reescreve a formulação a partir de duas fontes:
\begin{itemize}
\item \textbf{[N]} as notas \emph{Formulação poroelástica não saturada: esqueleto deformável, água e ar}
(\path{docs/biot_nao_saturado/biot_nao_saturado_agua_ar3.tex}), citadas por seção, \nota{§3};
\item \textbf{[A]} o artigo de Lira Cecílio (2026), \emph{Analytical return mapping for Mohr--Coulomb plasticity in
rotated Haigh--Westergaard space with consistent tangent operator}, J.\ Eng.\ Math.\ 157:10, citado por equação,
\art{Eq.~57}.
\end{itemize}
Ao lado de cada bloco está o trecho do NeoPZ que o resolve, \loc{arquivo}{linhas}. Onde o código difere das fontes
há uma marca \disc{D$k$}, explicada na tabela da Seção~\ref{sec:disc}.

\textbf{O que foi verificado.}
\begin{itemize}
\item \textbf{Mohr--Coulomb.} O programa do talude usa o passo \code{TPZPlasticStepVoigt} com o critério
\code{TPZYCMohrCoulombPV2}, a classe do artigo. As expressões do código são iguais às Tabelas~4--6 e às Eqs.~(57)--(58) do artigo. A região escolhida pelo
código coincide com a projeção de ponto mais próximo na norma de energia em 20\,000 estados aleatórios. O operador
tangente passa no teste de Taylor com ordem~2 (Seção~\ref{sec:mcverif}).
\item \textbf{Balanços de massa.} Com $\chi=S_w$, $\alpha=1$ e grãos incompressíveis, os resíduos do código são as
equações das notas discretizadas por Euler implícito. A diferença de fundo está no armazenamento: o código usa a forma
secante $n(S_w^{n+1}-S_w^n)$, que conserva a massa de água, e as notas usam a forma em taxa $C_{\pi\beta}\Delta p_\beta$.
\end{itemize}

\textbf{Discrepâncias que mudam resultados.}
\begin{itemize}
\item \disc{D1} O padrão do código é $\chi=1$ constante; as notas usam $\chi=S_w$.
\item \disc{D9} O termo de compressibilidade do ar usa $S_a^{n+1}$ no lugar de $S_a^{n}$.
\item \disc{D14} As rodadas com $\psi<\phi$ do manuscrito usaram a classe iterativa antiga \code{TPZYCMohrCoulombPV}, e
não o PV2 do artigo, que é só associado.
\item \disc{D3--D5} Porosidade, permeabilidade intrínseca e peso específico são constantes no código.
\end{itemize}
Três decisões ficam para você antes da Fase~2 (Seção~\ref{sec:decisoes}). O plano de testes está na
Seção~\ref{sec:plano}.

\tableofcontents

%======================================================================
\section{Hipóteses, variáveis primárias e convenções}\label{sec:hip}
%======================================================================

As notas \nota{§1} adotam pequenas deformações, regime quase estático, temperatura constante e duas fases fluidas
imiscíveis (água e ar) sem mudança de fase. O fluxo segue Darcy generalizado com permeabilidades relativas, e a
contribuição das pressões à tensão total segue Bishop, como escolha constitutiva. As variáveis primárias são
\begin{equation}
\boxed{\mathbf u(\mathbf x,t),\qquad p_w(\mathbf x,t),\qquad p_a(\mathbf x,t),}\qquad
\mathbf z=(\mathbf u,p_w,p_a).
\end{equation}

O que o código acrescenta ou fixa:
\begin{itemize}
\item \textbf{Esqueleto elastoplástico.} As notas tratam o esqueleto como elástico \nota{§2, §4}. No código a tensão
efetiva vem do Mohr--Coulomb associado do artigo (Seção~\ref{sec:mc}).
\item \textbf{Estado plano de deformação}, com $\varepsilon_{xz}=\varepsilon_{yz}=\varepsilon_{zz}=0$; o modelo constitutivo
trabalha com o tensor 3D e devolve $\sigma_{zz}$. \loc{TPZMatUnsaturatedUPwPa.cpp}{185--199}
\item \textbf{Unidades e sinais.} kN, m, s; pressões em kPa. Tensões positivas em tração; $p_w$ e $p_a$ manométricas e
positivas em compressão. $\mathbf g=(0,-1)$ é o vetor unitário da gravidade, e o peso entra por pesos específicos
$\gamma$ em kN/m$^3$. \loc{TPZMatUnsaturatedUPwPa.h}{57--59, 274}
\item \textbf{Deformação em Voigt com distorção de engenharia}, na ordem $xx,xy,xz,yy,yz,zz$, igual ao mapa de
deformação do artigo \art{Eq.~22} e à matriz elástica \art{Eq.~26}. \loc{TPZElasticResponse.cpp}{68--87, 127--140}
\end{itemize}

%======================================================================
\section{Cinemática e constantes elásticas}\label{sec:cinem}
%======================================================================

\nota{§2}, com $K$ o módulo volumétrico drenado do esqueleto (o $K$ do artigo):
\begin{equation}
\boxed{\eps(\mathbf u)=\tfrac12\left[\nabla\mathbf u+(\nabla\mathbf u)^{\mathsf T}\right],\qquad
\varepsilon_v=\Div\mathbf u,}
\end{equation}
\begin{equation}
G=\frac{E}{2(1+\nu)},\qquad \lambda=\frac{E\nu}{(1+\nu)(1-2\nu)},\qquad K=\frac{E}{3(1-2\nu)},
\end{equation}
\begin{equation}
\boxed{\alpha=1-\frac{K}{K_s},\qquad \eta=\frac{\alpha-n}{K_s}.}
\end{equation}
As notas admitem $\alpha(n(\mathbf z))$. O código usa $\alpha$ e $K_s$ constantes e, no talude, $\alpha=1$ com grãos
incompressíveis ($K_s\to\infty$, logo $\eta=0$). \loc{UnsatSlope.h}{184} \loc{TPZMatUnsaturatedUPwPa.cpp}{302--303}
\disc{D2}

%======================================================================
\section{Sucção, saturação e pressão equivalente}\label{sec:succao}
%======================================================================

\nota{§3}:
\begin{equation}
\boxed{p_c=p_a-p_w,\qquad S_w=S_w(p_c),\qquad S_a=1-S_w,\qquad
\Gamma=\frac{\partial S_w}{\partial p_w}=-\frac{\partial S_w}{\partial p_a}=-\frac{\dd S_w}{\dd p_c}\ge0.}
\end{equation}
Para $p_c\le0$, $S_w=1$ e $\Gamma=0$. \loc{TPZUnsaturatedHydraulics.h}{78--83}

\textbf{Curvas de retenção do código} \loc{TPZUnsaturatedHydraulics.h}{78--99}:
\begin{itemize}
\item van Genuchten, igual às notas \nota{§3}: $S_e=[1+(a_{\mathrm{VG}}p_c)^{b_{\mathrm{VG}}}]^{-m_{\mathrm{VG}}}$,
$m_{\mathrm{VG}}=1-1/b_{\mathrm{VG}}$, $S_w=S_{wr}+(1-S_{wr})S_e$, com o mesmo $\Gamma$ das notas;
\item Fredlund--Xing com fator de correção $C(\psi)=1$: $S_w=[\ln(e+(p_c/a)^n)]^{-m}$;
\item Liakopoulos: $S_w=1-c_L\,p_c^{\,b_L}$, limitada a $S_{wr}$.
\end{itemize}
O padrão do programa do talude é Fredlund--Xing, e não van Genuchten. \loc{UnsatSlope.h}{121--131} \disc{D7}

\textbf{Pressão equivalente.} As notas fixam $\chi=S_w$ \nota{§1, §3}:
\begin{equation}
\boxed{p_B=S_wp_w+S_ap_a=p_a-S_wp_c.}
\label{eq:pB}
\end{equation}
O código admite três leis \loc{TPZMatUnsaturatedUPwPa.cpp}{97--118}:
\begin{equation}
p_B=\chi p_w+(1-\chi)p_a,\qquad
\chi=\begin{cases}
S_w & \text{(notas)},\\
S_e=(S_w-S_{wr})/(1-S_{wr}),\\
\chi_0 \text{ se } p_w\le p_a,\quad 1 \text{ se } p_w>p_a & \text{(padrão, com } \chi_0=1).
\end{cases}
\end{equation}
Com o padrão, $p_B=p_w$ em toda a zona não saturada ($\phi^b=\phi'$).
\loc{TPZMatUnsaturatedUPwPa.h}{275--276} \loc{UnsatSlope.h}{114--115} \disc{D1}

\textbf{Coeficientes de Bishop tangentes} \nota{§13.4}, usados no Jacobiano:
\begin{equation}
\boxed{\chi_w=\frac{\partial p_B}{\partial p_w}=S_w-\Gamma p_c,\qquad
\chi_a=\frac{\partial p_B}{\partial p_a}=S_a+\Gamma p_c,\qquad \chi_w+\chi_a=1.}
\end{equation}
O código calcula exatamente estes valores com $\chi=S_w$. Com $\chi=S_e$ usa
$\chi_w=S_e-(\dd S_e/\dd S_w)\,\Gamma p_c$. Com $\chi$ constante usa $\chi_w=\chi$, que salta de $\chi_0$ para 1 em
$p_c=0$ quando $\chi_0\neq1$. \loc{TPZMatUnsaturatedUPwPa.cpp}{100--116}

%======================================================================
\section{Equilíbrio mecânico}\label{sec:equil}
%======================================================================

\nota{§4}, tração positiva:
\begin{equation}
\boxed{\Div\sig+\rho_b\mathbf g=\mathbf0,\qquad \sig=\sig'-\alpha\,p_B\,\Id.}
\end{equation}
As notas escrevem a forma elástica incremental $\sig=\sig_0+2G\eps+\lambda\varepsilon_v\Id-\alpha(p_B-p_{B0})\Id$. No
código, $\sig'$ é a saída do Mohr--Coulomb (Seção~\ref{sec:mc}) para a deformação total medida a partir do estado
inicial, e $p_B$ entra com seu valor total. As duas escritas coincidem porque o estado inicial é calculado pelo próprio
modelo (Seção~\ref{sec:cc}). O carregamento geostático equilibra o peso com as pressões iniciais; depois os
deslocamentos são zerados sem alterar $\sig'$ ($\eps^p\leftarrow\eps^p-\eps$, $\eps\leftarrow\mathbf0$). Logo
$\sig_0=\sig'_0-\alpha p_{B0}\Id$. \loc{TPZMatUnsaturatedUPwPa.cpp}{430--438}

\textbf{Resíduo do equilíbrio no código} \loc{TPZMatUnsaturatedUPwPa.cpp}{254--267}:
\begin{equation}
\mathbf R_u=\int_\Omega\mathbf B_u^{\mathsf T}\sig'\,\dd\Omega
-\lambda_p\,\alpha\int_\Omega\mathbf B_u^{\mathsf T}\mathbf m\,p_B\,\dd\Omega
-\lambda_g\int_\Omega\mathbf N_u^{\mathsf T}\gamma\,\mathbf g\,\dd\Omega-\mathbf f_t .
\label{eq:Ru}
\end{equation}
Os fatores $\lambda_g$ (peso) e $\lambda_p$ (pressão equivalente) não existem nas notas:
\begin{itemize}
\item no transiente, $\lambda_g=\lambda_p=1$;
\item no carregamento geostático, $\lambda_p=\lambda_g$ cresce de 0 a 1 junto com o peso;
\item no fator de segurança por aumento de gravidade, $\lambda_p=\lambda_g$ onde $p_B>0$ e $\lambda_p=1$ onde $p_B<0$: a
sucção não é escalada (Seção~\ref{sec:fs}). \loc{TPZMatUnsaturatedUPwPa.cpp}{232--233}
\end{itemize}

\textbf{Peso.} As notas usam a massa específica da mistura, $\rho_b=(1-n)\rho_s+n(S_w\rho_w+S_a\rho_a)$ \nota{§4}. O
código usa por padrão $\gamma$ constante, 20~kN/m$^3$. Como opção, usa $\gamma=(1-n)\gamma_s+nS_w\gamma_w$, sem o
peso do ar e com $n$ constante. \loc{TPZMatUnsaturatedUPwPa.cpp}{120--129} \disc{D5}

%======================================================================
\section{Lei de Darcy e permeabilidades}\label{sec:darcy}
%======================================================================

\nota{§5}:
\begin{equation}
\boxed{\mathbf q_\pi=-\lambda_\pi\,(\nabla p_\pi-\rho_\pi\mathbf g),\qquad
\lambda_\pi=\frac{\kappa\,k_{r\pi}}{\mu_\pi},\qquad \pi\in\{w,a\}.}
\end{equation}
O código escreve as mobilidades com a condutividade saturada $k_s$ (m/s) e o peso específico da água
\loc{TPZMatUnsaturatedUPwPa.cpp}{313--322}:
\begin{equation}
\mathbf q_w=-\frac{k_{rw}}{\gamma_w}\,\mathbf K(\nabla p_w-\gamma_w\mathbf g),\qquad
\mathbf q_a=-\frac{\mu_w}{\mu_a}\,\frac{k_{ra}}{\gamma_w}\,\mathbf K(\nabla p_a-\gamma_a\mathbf g),\qquad
\mathbf K=k_s\Id .
\label{eq:darcycod}
\end{equation}
É a forma das notas com $\kappa=k_s\mu_w/\gamma_w$ e $\rho_\pi\mathbf g=\gamma_\pi\mathbf g$. O tensor anisotrópico
$\mathbf K=\mathbf R(\theta)\operatorname{diag}(k_1,k_2)\mathbf R(\theta)^{\mathsf T}$ é uma opção que as notas não têm.
\loc{TPZMatUnsaturatedUPwPa.h}{159--170}

Diferenças: as notas fazem $\kappa(n)$ variar com a porosidade (Kozeny--Carman), e no código $\mathbf K$ é constante
\disc{D4}; o peso específico do ar $\gamma_a$ é constante (1,2~kg/m$^3$) e não $\rho_a(p_a)\,g$ \disc{D6}.

\textbf{Permeabilidades relativas} \loc{TPZUnsaturatedHydraulics.h}{108--162}:
\begin{center}\small
\begin{tabular}{@{}p{0.09\linewidth}p{0.33\linewidth}p{0.52\linewidth}@{}}
\toprule
& notas \nota{§5, §13.9} & código (\textbf{padrão do talude} em negrito)\\
\midrule
água & Mualem: $S_e^{1/2}\{1-[1-S_e^{1/m}]^{m}\}^2$ & Mualem, igual às notas; \textbf{potência $S_e^{4}$}; Liakopoulos\\
ar & Mualem: $(1-S_e)^{1/2}[1-S_e^{1/m}]^{2m}$ & Mualem, igual às notas; \textbf{Brooks--Corey
$(1-S_e)^2(1-S_e^{(2+\lambda)/\lambda})$, $\lambda=3$}\\
limites & $k_{ra}\ge k_{ra}^{\min}$ como opção & $k_{rw}\ge10^{-12}$ e $k_{ra}\ge10^{-4}$, com derivada nula no limite\\
\bottomrule
\end{tabular}
\end{center}
\disc{D7}

%======================================================================
\section{Compressibilidade e massa específica dos fluidos}\label{sec:fluidos}
%======================================================================

\nota{§6}: $\rho_w=\rho_{w0}\exp[(p_w-p_{w0})/K_w]$; para o ar ideal isotérmico, $\rho_a\propto\pabs$ e $K_a=\pabs$, com
$\pabs=p_a+p_{\mathrm{atm}}$. O termo de transporte normalizado se expande como \nota{§13.2}
\begin{equation}
\frac{1}{\rho_\pi}\Div(\rho_\pi\mathbf q_\pi)=\Div\mathbf q_\pi+c_\pi\,\mathbf q_\pi\cdot\nabla p_\pi,\qquad
c_w=\frac1{K_w},\qquad c_a=\frac1{\pabs}.
\end{equation}
No código \loc{TPZMatUnsaturatedUPwPa.cpp}{320--330}:
\begin{itemize}
\item a compressibilidade da água entra só pelo armazenamento $nS_w\Delta p_w/K_w$; o termo $c_w\,\mathbf q_w\cdot\nabla p_w$
não existe, como as próprias notas admitem \nota{§13.4};
\item o ar entra pelo armazenamento $nS_a\Delta p_a/\pabs$ e pelo termo $c_a\,\mathbf q_a\cdot\nabla p_a$, completo e com
derivada no Jacobiano;
\item $K_w=2{,}2\times10^6$~kPa e $p_{\mathrm{atm}}=101{,}325$~kPa. \loc{TPZMatUnsaturatedUPwPa.h}{269--271}
\end{itemize}

%======================================================================
\section{Armazenamento e acoplamento entre as fases}\label{sec:armaz}
%======================================================================

\nota{§7}, com $D_w=n-\eta S_wp_c$ e $D_a=-n-\eta S_ap_c$:
\begin{equation}
\boxed{\begin{aligned}
B_w&=\alpha S_w, & B_a&=\alpha S_a,\\
C_{ww}&=\eta S_w^2+\frac{nS_w}{K_w}+D_w\Gamma, & C_{wa}&=\eta S_wS_a-D_w\Gamma,\\
C_{aw}&=\eta S_aS_w+D_a\Gamma, & C_{aa}&=\eta S_a^2+\frac{nS_a}{\pabs}-D_a\Gamma,
\end{aligned}}
\end{equation}
que vêm de $\dot\zeta_i=S_i(\alpha\dot\varepsilon_v+\eta\dot p_B)+n\dot S_i+(nS_i/K_i)\dot p_i$, $i\in\{w,a\}$.

O código usa os mesmos termos, com $\eta\,\Delta p_B$ calculado como diferença secante $p_B^{n+1}-p_B^n$, que já
contém a parcela $(p_w-p_a)\Delta S_w$. \loc{TPZMatUnsaturatedUPwPa.cpp}{301--312} Com os valores do talude
($\alpha=1$, $\eta=0$, logo $D_w=n$ e $D_a=-n$):
\begin{equation}
B_w=S_w,\quad B_a=S_a,\quad
C_{ww}=\frac{nS_w}{K_w}+n\Gamma,\quad C_{wa}=C_{aw}=-n\Gamma,\quad C_{aa}=\frac{nS_a}{\pabs}+n\Gamma.
\label{eq:Creduzido}
\end{equation}
O acoplamento $B_w=\alpha S_w$ do balanço de massa não depende do $\chi$ escolhido para a tensão. Com $\chi=S_w$ é o par
derivado nas notas. Com o padrão $\chi=1$, o equilíbrio usa $\chi_w=1$ e a massa usa $S_w$ \disc{D1}.

%======================================================================
\section{Balanços de massa e sistema acoplado}\label{sec:sistema}
%======================================================================

\nota{§8, §10}:
\begin{equation}
\boxed{\begin{aligned}
\Div\sig+\rho_b\mathbf g&=\mathbf0,\\
B_w\dot\varepsilon_v+C_{ww}\dot p_w+C_{wa}\dot p_a+\Div\mathbf q_w+c_w\,\mathbf q_w\cdot\nabla p_w&=s_w,\\
B_a\dot\varepsilon_v+C_{aw}\dot p_w+C_{aa}\dot p_a+\Div\mathbf q_a+c_a\,\mathbf q_a\cdot\nabla p_a&=s_a.
\end{aligned}}
\label{eq:sistema}
\end{equation}

\textbf{Porosidade.} As notas fecham $n=n_0+(\alpha_0-n_0)\varepsilon_v+(\alpha_0-n_0)(p_B-p_{B0})/K_s$ em primeira
ordem \nota{§9}. O código mantém $n=n_0$ em todos os coeficientes; a variação do volume dos poros entra só pelo termo
$B_\pi\dot\varepsilon_v$. \loc{TPZMatUnsaturatedUPwPa.cpp}{300} \disc{D3}

\textbf{Limite saturado} \nota{§12}. Com $S_w=1$ e $\Gamma=0$, a equação da água fica
$\alpha\dot\varepsilon_v+(n/K_w)\dot p_w+\Div\mathbf q_w=0$ (com $\eta=0$), que é Biot saturado. O código reproduz
esse limite pela Eq.~\eqref{eq:Creduzido}.

\textbf{Desaparecimento do ar.} Quando $S_a\to0$, $C_{aa}$ e $\mathbf H_a$ tendem a zero \nota{§13.9}. O código adota a
primeira saída das notas, $k_{ra}\ge k_{ra}^{\min}=10^{-4}$, e mantém $p_a$ como incógnita em todo o domínio. Na zona
saturada a equação do ar vira uma difusão com mobilidade mínima. \loc{TPZUnsaturatedHydraulics.h}{155--161}

%======================================================================
\section{Condições de contorno e estado inicial do talude}\label{sec:cc}
%======================================================================

As notas dão a forma geral (Dirichlet e Neumann) \nota{§13.1}. No exemplo da barragem, a face de percolação é a
complementaridade $p_w\le0$, $\mathbf q_w\cdot\mathbf n\ge0$, $p_w\,(\mathbf q_w\cdot\mathbf n)=0$ \nota{§16}. O programa
do talude usa \loc{main.cpp}{355--401}:
\begin{itemize}
\item \textbf{Superfície (crista, face e pé): chuva unilateral por conjunto ativo.} Cada vértice recebe a chuva $I$ na
projeção horizontal, metade de cada aresta vizinha. Quando $p_w$ chega a zero o vértice passa a ter $p_w=0$ prescrita,
e volta a receber fluxo quando a vazão que entra por ele excede a chuva. É a complementaridade das notas com a chuva no
lugar do fluxo nulo: $p_w\le0$, $q_{\mathrm{in}}\le I$, $p_w\,(I-q_{\mathrm{in}})=0$. A mesma estrutura está nas
Eqs.~(23)--(25) de Pedroso (2015, p.~6). Até 12 passagens por passo. \loc{main.cpp}{523--610}
\item \textbf{Ar.} Modelo ativo: $p_a=0$ em todos os vértices da superfície, inclusive os alagados. Modelo passivo:
$p_a=0$ em todo o domínio. Base e laterais são impermeáveis ao ar, salvo opção. \loc{main.cpp}{391--399}
\item \textbf{Água.} Pressão hidrostática prescrita nas laterais abaixo do nível d'água; base impermeável.
\loc{main.cpp}{378--390}
\item \textbf{Deslocamentos.} Base fixa e laterais com rolete. \loc{UnsatSlope.h}{525--532}
\item \textbf{Estado inicial.} $p_w$ hidrostática abaixo do nível d'água e sucção limitada a $s_{\max}=75$~kPa acima;
$p_a=0$. O peso é aplicado em quatro incrementos com as pressões congeladas; se não convergir com a resistência real, é
aplicado com resistência aumentada ($F_0=0{,}6$) e a resistência é reduzida de volta. Depois $\mathbf u=\mathbf0$ sem
mudar $\sig'$. \loc{main.cpp}{404--410, 813--851}
\end{itemize}
O campo inicial de pressões não é estacionário: a sucção limitada a 75~kPa não é solução do fluxo sem chuva \disc{D16}.
Com $p_a=0$ fixo na superfície alagada, o aprisionamento do ar depende de $k_{ra}\to k_{ra}^{\min}$ na camada saturada
\disc{D17}.

%======================================================================
\section{Mohr--Coulomb associado na forma do artigo}\label{sec:mc}
%======================================================================

Esta seção segue o artigo \art{Seções 2--4}. A projeção está em \path{Material/Plasticity/TPZYCMohrCoulombPV2.cpp}; o
passo plástico e a tangente, em \path{Material/Plasticity/TPZPlasticStepVoigt.cpp}.

\subsection{Critério e setor ordenado}
Com $\sigma_1\ge\sigma_2\ge\sigma_3$ e tração positiva, só três dos seis planos podem estar ativos
\art{Eqs.~44, 45, 49}:
\begin{equation}
\begin{aligned}
\Phi_1&=(\sigma_1-\sigma_3)+(\sigma_1+\sigma_3)\sn-2c\cs,\\
\Phi_2&=(\sigma_2-\sigma_3)+(\sigma_2+\sigma_3)\sn-2c\cs,\\
\Phi_6&=(\sigma_1-\sigma_2)+(\sigma_1+\sigma_2)\sn-2c\cs.
\end{aligned}
\end{equation}
\loc{TPZYCMohrCoulombPV2.cpp}{46--53}. As tensões principais chegam ordenadas: \code{TPZTensor} calcula os autovalores
pela fórmula trigonométrica e para se a ordem falhar. \loc{Matrix/TPZTensor.h}{1657--1674}

\subsection{Espaços de Haigh--Westergaard (HW) e HW rotacionado (RHW)}
\art{Eqs.~1--14}: $\xi=I_1/\sqrt3$, $\rho=\sqrt{2J_2}$ e
$\beta=\tfrac13\arccos\!\big(\tfrac{3\sqrt3}{2}J_3/J_2^{3/2}\big)\in[0,\pi/3]$, com
\begin{equation}
\sig^*=\begin{bmatrix}\xi\\ \rho\cos\beta\\ \rho\sin\beta\end{bmatrix}=\mathbf R\,\sig,\qquad
\mathbf R=\begin{bmatrix}\tfrac1{\sqrt3}&\tfrac1{\sqrt3}&\tfrac1{\sqrt3}\\[2pt]
\sqrt{\tfrac23}&-\tfrac1{\sqrt6}&-\tfrac1{\sqrt6}\\[2pt]
0&\tfrac1{\sqrt2}&-\tfrac1{\sqrt2}\end{bmatrix},
\end{equation}
\begin{equation}
\mathbf D_{RHW}=\mathbf R\,\mathbf D_{HW}\mathbf R^{\mathsf T}=\operatorname{diag}(3K,\,2G,\,2G).
\end{equation}
Nesta base a resposta elástica é diagonal e a parte hidrostática se separa da desviadora.

\subsection{Projeção de ponto mais próximo na norma de energia}
\art{Eqs.~15--18}: o estado projetado é o ponto admissível mais próximo da tentativa elástica na norma de energia
complementar,
\begin{equation}
\sig^*_{\mathrm{proj}}=\arg\min_{\Phi(\sig^*)\le0}\ \big\|\sig^*_{\mathrm{trial}}-\sig^*\big\|_{RHW},\qquad
\|\sig^*\|_{RHW}^2=\sig^{*\mathsf T}\mathbf D_{RHW}^{-1}\sig^*.
\end{equation}
Na face principal, $\Phi_1=0$ elimina $\xi$ \art{Eq.~52}:
\begin{equation}
\xi(\rho,\beta)=\frac{12c\cs-\sqrt6\,\rho(3+\sn)\cos\beta+3\sqrt2\,\rho(\sn-1)\sin\beta}{4\sqrt3\,\sn},
\end{equation}
e a projeção é o ponto estacionário da distância em $(\rho,\beta)$ \art{Eqs.~53--55}. Como a superfície é poliédrica no
setor, o mínimo está no interior da face principal ou em uma das arestas, o que permite a atualização em forma fechada.

\subsection{Atualização em forma fechada e escolha da região}\label{sec:mccand}
\art{Seção~4.4}, implementada em \code{ProjectSigma} \loc{TPZYCMohrCoulombPV2.cpp}{265--316}:
\begin{enumerate}
\item \textbf{Elástico} se $\Phi_1(\sig_{\mathrm{trial}})\le0$ \art{Eq.~56}. (linhas 274--281)
\item \textbf{Face principal} \art{Eqs.~57--58}:
\begin{equation}
\begin{gathered}
\sigma_i^{\mathrm{proj}}=\frac{(\sigma_1+\sigma_3)A_{i1}+(\sigma_1-\sigma_3)A_{i2}+(\sigma_1+\sigma_2)A_{i3}+A_{i4}}{\mathcal D},\\
\Big[\frac{\partial\sigma_i^{\mathrm{proj}}}{\partial\sigma_1},\ \frac{\partial\sigma_i^{\mathrm{proj}}}{\partial\sigma_2},\
\frac{\partial\sigma_i^{\mathrm{proj}}}{\partial\sigma_3}\Big]=\frac1{\mathcal D}
\big[A_{i1}+A_{i2}+A_{i3},\ \ A_{i3},\ \ A_{i1}-A_{i2}\big],
\end{gathered}
\label{eq:mcproj}
\end{equation}
com os coeficientes da Tabela~\ref{tab:A}. É aceita se preserva a ordem
$\sigma_1^{\mathrm{proj}}\ge\sigma_2^{\mathrm{proj}}-\epsilon_{\mathrm{ord}}$ e
$\sigma_2^{\mathrm{proj}}\ge\sigma_3^{\mathrm{proj}}-\epsilon_{\mathrm{ord}}$ \art{Eq.~59}. (linhas 95--138 e 284--290)
\item \textbf{Aresta}, se a ordem falha. Direita ($\beta=0$, $\sigma_2=\sigma_3$) se
$\chi_{\mathrm{MC}}=(1-\sn)\sigma_1-2\sigma_2+(1+\sn)\sigma_3>0$; esquerda ($\beta=\pi/3$, $\sigma_1=\sigma_2$) caso
contrário \art{Eq.~60}. Usa as mesmas fórmulas \eqref{eq:mcproj} com os coeficientes de aresta da
Tabela~\ref{tab:A}. (linhas 141--237 e 291--311)
\item \textbf{Ápice} se a aresta também viola a ordem: $\sigma_i^{\mathrm{proj}}=c\cot\phi$ e
$\partial\sig^{\mathrm{proj}}/\partial\sig=\mathbf0$ \art{Eq.~61}. (linhas 239--257 e 312--314)
\end{enumerate}
No código, $\epsilon_{\mathrm{ord}}=10^{-12}$ absoluto (função \code{IsZero}). \loc{Common/pzreal.h}{737--742}

\begin{table}[htbp]\centering\footnotesize\setlength{\tabcolsep}{3.5pt}
\caption{Coeficientes $A_{ij}$ e denominadores $\mathcal D$ das Tabelas~4--6 do artigo, com $s=\sn$ e $c_\phi=\cs$.
Transcritos do artigo e conferidos contra o código (Seção~\ref{sec:mcverif}).}\label{tab:A}
\renewcommand{\arraystretch}{1.3}
\begin{tabular}{@{}llllll@{}}
\toprule
região & $i$ & $A_{i1}$ & $A_{i2}$ & $A_{i3}$ & $A_{i4}$\\
\midrule
face & 1 & $-3G(s-1)$ & $(G+3K)(s-1)s$ & $0$ & $2cc_\phi[3Ks+G(3+s)]$\\
principal & 2 & $G(s^2-3)-6Ks^2$ & $-3Ks(1+s)-G[3+(s-2)s]$ & $6G+2(G+3K)s^2$ & $(6K-4G)\,c\,c_\phi s$\\
 & 3 & $3G(1+s)$ & $-(G+3K)s(1+s)$ & $0$ & $2cc_\phi[G(s-3)+3Ks]$\\
\multicolumn{6}{@{}l}{\quad $\mathcal D_f=2[3G+(G+3K)s^2]$}\\
\midrule
aresta & 1 & $-\tfrac12(s-1)a$ & $\tfrac12(s-1)b$ & $-(s-1)d$ & $4cc_\phi d$\\
direita & 2 e 3 & $\tfrac12(1+s)a$ & $-\tfrac12(1+s)b$ & $(1+s)d$ & $-2cc_\phi a$\\
\multicolumn{6}{@{}l}{\quad $\mathcal D_d=9G+6Gs+Gs^2+12Ks^2$, \ $a=G(3+s)-6Ks$, \ $b=G(3+s)+12Ks$, \ $d=G(3+s)+3Ks$}\\
\midrule
aresta & 1 e 2 & $\tfrac12(s-1)a'$ & $-\tfrac12(s-1)a'$ & $(s-1)d'$ & $-2cc_\phi a'$\\
esquerda & 3 & $-\tfrac12(1+s)a'$ & $\tfrac12(1+s)a'$ & $-(1+s)d'$ & $4cc_\phi d'$\\
\multicolumn{6}{@{}l}{\quad $\mathcal D_e=9G-6Gs+Gs^2+12Ks^2$, \ $a'=G(s-3)-6Ks$, \ $d'=G(s-3)+3Ks$}\\
\bottomrule
\end{tabular}
\end{table}

\subsection{Reconstrução da tensão}
\art{Eq.~62}: $\sig^{\mathrm{proj}}=\sum_i\sigma_i^{\mathrm{proj}}\,\mathbf v_i\mathbf v_i^{\mathsf T}$, com os autovetores
$\mathbf v_i$ da tentativa. O código soma só a correção plástica,
$\sig=\sig_{\mathrm{trial}}+\sum_i(\sigma_i^{\mathrm{proj}}-\sigma_i^{\mathrm{trial}})\,\mathbf v_i\mathbf v_i^{\mathsf T}$,
que é a mesma tensão com menos arredondamento na parte elástica. \loc{TPZPlasticStepVoigt.cpp}{193--201}

\subsection{Operador tangente consistente}\label{sec:mctan}
\art{Eqs.~25--37 e 43}:
\begin{equation}
\mathbb D=\mathbb A+\mathbb R_V,\qquad
\mathbb A=\Big(\sum_{i,j}D_{ij}\,\underline{\mathbf v}^{\sigma}_{ii}\,(\underline{\mathbf v}^{\varepsilon}_{jj})^{\mathsf T}\Big)\,\mathbb C,\qquad
\mathbb R_V=2\sum_{i<j}\kappa_{ij}\,\underline{\mathbf s}_{ij}\,\underline{\mathbf s}_{ij}^{\mathsf T},\qquad
D_{ij}=\frac{\partial\sigma_i^{\mathrm{proj}}}{\partial\sigma_j},
\end{equation}
\begin{equation}
\kappa_{ij}=\frac{\sigma_i^{\mathrm{proj}}-\sigma_j^{\mathrm{proj}}}{\varepsilon_i-\varepsilon_j},
\qquad\text{ou}\qquad
\kappa_{ij}=G\,(D_{ii}-D_{ij}-D_{ji}+D_{jj})\quad\text{se } |\varepsilon_i-\varepsilon_j|<\epsilon_\kappa .
\end{equation}
Aqui $\underline{\mathbf v}^{\sigma}_{ii}=\mathsf V_\sigma[\mathbf v_i\otimes\mathbf v_i]$,
$\underline{\mathbf v}^{\varepsilon}_{jj}=\mathsf V_\varepsilon[\mathbf v_j\otimes\mathbf v_j]$,
$\underline{\mathbf s}_{ij}=\mathsf V_\sigma[\tfrac12(\mathbf v_i\mathbf v_j^{\mathsf T}+\mathbf v_j\mathbf v_i^{\mathsf T})]$, e
$\varepsilon_i$ são as deformações elásticas principais da tentativa. O código monta $\mathbb A$ e $\mathbb R_V$ coluna a
coluna, como a Listing~1 do artigo, com $\epsilon_\kappa=10^{-12}$; a Listing usa $10^{-15}$. No ápice, $\mathbb D=\mathbf0$.
\loc{TPZPlasticStepVoigt.cpp}{214--282}

\subsection{Redução de resistência}
Não está no artigo. O código aplica, a uma cópia do critério, $c\to c/F$, $\tan\phi\to\tan\phi/F$ e
$\tan\psi\to\tan\psi/F$. Como $\psi=\phi$, a regra continua associada para todo $F$.
\loc{TPZYCMohrCoulombPV2.cpp}{75--81} \loc{TPZPlasticStepVoigt.cpp}{78--90}

\subsection{Regra não associada}
O PV2 é associado: o código para (\code{DebugStop}) se $\psi\neq\phi$, e o programa do talude recusa $\psi\neq\phi$.
\loc{TPZYCMohrCoulombPV2.cpp}{273} \loc{UnsatSlope.h}{143--146} As rodadas com $\psi<\phi$ do manuscrito foram feitas
com o executável \code{SlopeRainfallPV}, que troca o modelo pela classe iterativa antiga
\code{TPZPlasticStepPV<TPZYCMohrCoulombPV>}. Nesse executável o FS por SRM é calculado com $\psi=\phi$.
\loc{UnsatSlope.h}{40--46} \loc{Projects/SlopeRainfall/CMakeLists.txt}{6--8} \loc{main.cpp}{447--458} \disc{D14}

\subsection{Verificação numérica do código contra o artigo}\label{sec:mcverif}

\textbf{(a) Expressões e escolha da região.} O script \path{verificacao/verifica_pv2.py} lê as expressões C++ de
\path{TPZYCMohrCoulombPV2.cpp} e as avalia em 20\,000 estados aleatórios, com $E\in[10^3,10^7]$~kPa,
$\nu\in[0;\,0{,}49]$, $\phi\in[5^\circ,45^\circ]$ e $c\in[1,1000]$~kPa. A referência independente é a projeção de ponto
mais próximo por enumeração de conjuntos ativos dos seis planos, com as condições de Karush--Kuhn--Tucker. Saída em
\path{verificacao/verifica_pv2.txt}:
\begin{center}\small
\begin{tabular}{@{}lr@{}}
\toprule
comparação & diferença máxima\\
\midrule
tensões do código $\times$ Tabelas~4--6 e Eq.~57 (relativa) & $1{,}7\times10^{-15}$\\
Jacobiano do código $\times$ Eq.~58 & $1{,}1\times10^{-15}$\\
Jacobiano do código $\times$ diferenças centrais & $9{,}3\times10^{-10}$\\
tensão escolhida pelo código $\times$ ponto mais próximo (relativa) & $7{,}1\times10^{-15}$\\
\bottomrule
\end{tabular}
\end{center}
Os estados caíram em todas as regiões: elástico 2660, face principal 6652, aresta esquerda 5431, aresta direita 2601,
ápice 2656.

\textbf{(b) Teste de Taylor} \art{Eqs.~63--65}, com o executável \code{SlopeMohrCoulomb check}; saída em
\path{verificacao/taylor_check.txt}:
\begin{center}\small
\begin{tabular}{@{}llll@{}}
\toprule
estado & região & ordem $p$ ou erro & $|\mathbb D-\mathbb D^{\mathsf T}|/|\mathbb D|$\\
\midrule
geral & face principal & 1,99991 & $2{,}3\times10^{-15}$\\
$\varepsilon_2=\varepsilon_3$ & aresta direita & 1,90078 & $2{,}4\times10^{-15}$\\
$\varepsilon_1=\varepsilon_2$ & aresta esquerda & 1,99994 & $2{,}5\times10^{-15}$\\
tração & ápice & $|E|=4{,}6\times10^{-13}$ (exato) & 0\\
geral & face principal & 1,99889 & $5{,}7\times10^{-15}$\\
pequeno & elástico & $|E|=4{,}8\times10^{-15}$ (exato) & 0\\
\bottomrule
\end{tabular}
\end{center}

\textbf{(c) Observações sobre o texto do artigo.} Não afetam o código. As três primeiras já constavam do README de
\path{Projects/SlopeMohrCoulomb}; a primeira e a quarta foram conferidas de novo para este documento.
\begin{itemize}
\item A Eq.~(54) vale 432 vezes o quadrado da distância; o minimizador não muda (razão 432,000000 em cinco estados,
\path{verificacao/eq54.txt}).
\item As Tabelas~4--6 imprimem ``$\sin\phi^2$'' onde deve ser $\sin^2\phi$.
\item A Eq.~(28) deveria ter $(\underline{\mathbf v}^{\varepsilon}_{jj})^{\mathsf T}$, e não
$(\underline{\mathbf v}^{\sigma}_{jj})^{\mathsf T}$, e é só a parte coaxial; a Eq.~(29), implementada, está correta.
\item O texto após a Eq.~(14) diz $\mathbf D_{RHW}=\mathbf R^{\mathsf T}\mathbf D_{HW}\mathbf R$; com
$\sig^*=\mathbf R\sig$ (Eq.~10) o correto é $\mathbf R\,\mathbf D_{HW}\mathbf R^{\mathsf T}$, que é a matriz diagonal da
Eq.~(14). A outra ordem não dá matriz diagonal.
\end{itemize}

%======================================================================
\section{Forma fraca}\label{sec:fraca}
%======================================================================

\nota{§13.2}, com funções-teste $\mathbf v$ (nula em $\Gamma_u$) e $w$ (nula em $\Gamma_p^\pi$):
\begin{equation}
\int_\Omega\eps(\mathbf v):\sig\,\dd\Omega=\int_\Omega\mathbf v\cdot\rho_b\mathbf g\,\dd\Omega
+\int_{\Gamma_t}\mathbf v\cdot\bar{\mathbf t}\,\dd\Gamma,
\end{equation}
\begin{equation}
\int_\Omega w\,[B_\pi\dot\varepsilon_v+C_{\pi w}\dot p_w+C_{\pi a}\dot p_a]\,\dd\Omega
-\int_\Omega\nabla w\cdot\mathbf q_\pi\,\dd\Omega
+\int_\Omega w\,c_\pi\,\mathbf q_\pi\cdot\nabla p_\pi\,\dd\Omega
=\int_\Omega w\,s_\pi\,\dd\Omega-\int_{\Gamma^\pi_q}w\,\bar q_\pi\,\dd\Gamma .
\end{equation}
O código usa esta forma com as equações de massa multiplicadas por $\Delta t$ (Seção~\ref{sec:temporal}). Duas
convenções mudam. Nas notas $\bar q_\pi>0$ é saída; no código o fluxo prescrito (\code{EFluxW}, \code{EFluxA}) é positivo
entrando. A chuva do talude entra como vazão nodal nas equações da água dos vértices da superfície.
\loc{TPZMatUnsaturatedUPwPa.cpp}{413--419} \loc{UnsatSlope.h}{365} \disc{D20}

%======================================================================
\section{Discretização espacial}\label{sec:espacial}
%======================================================================

As notas usam interpolações de mesma ordem (Q4 ou T3 para todos os campos) \nota{§13.3, §14} e alertam que elas violam
a condição inf-sup no limite não drenado, recomendando Taylor--Hood \nota{§13.9}. O código usa Taylor--Hood em
triângulos: $\mathbf u$ quadrático (T6), $p_w$ e $p_a$ lineares (T3). \loc{UnsatSlope.h}{199--203} \disc{D12}

\textbf{Funções de forma}, em coordenadas de área $L_1,L_2,L_3$ ($L_1=1-\xi-\eta$, $L_2=\xi$, $L_3=\eta$):
\begin{equation}
\begin{aligned}
&\text{T6 }(\mathbf u): && N_i=L_i(2L_i-1)\ \ (i=1,2,3),\qquad N_4=4L_1L_2,\qquad N_5=4L_2L_3,\qquad N_6=4L_3L_1,\\
&\text{T3 }(p_w,p_a): && N_i=L_i=\frac{a_i+b_ix+c_iy}{2A}\ \ (i=1,2,3).
\end{aligned}
\end{equation}
Os coeficientes $a_i,b_i,c_i$ e a área $A$ do T3 são os das notas \nota{§14}. Para o T6, as derivadas $N_{i,x}$ e
$N_{i,y}$ são lineares no elemento e se obtêm de $\partial L_k/\partial x=b_k/2A$ e $\partial L_k/\partial y=c_k/2A$ (lados
retos). As matrizes têm a estrutura geral das notas \nota{§14}:
\begin{equation}
\mathbf N_u=\begin{bmatrix}N_1&0&\cdots&N_6&0\\0&N_1&\cdots&0&N_6\end{bmatrix}_{2\times12},\quad
\mathbf B_u=\begin{bmatrix}N_{1,x}&0&\cdots&N_{6,x}&0\\0&N_{1,y}&\cdots&0&N_{6,y}\\N_{1,y}&N_{1,x}&\cdots&N_{6,y}&N_{6,x}\end{bmatrix}_{3\times12},
\end{equation}
\begin{equation}
\mathbf N_p=\begin{bmatrix}N_1&N_2&N_3\end{bmatrix},\qquad
\mathbf B_p=\frac1{2A}\begin{bmatrix}b_1&b_2&b_3\\c_1&c_2&c_3\end{bmatrix},\qquad
\mathbf m^{\mathsf T}\mathbf B_u\mathbf U^e=\varepsilon_v .
\end{equation}
No código, $\mathbf B_u$ tem seis linhas na ordem Voigt 3D, com as linhas $xz$, $yz$ e $zz$ nulas.
\loc{TPZMatUnsaturatedUPwPa.cpp}{268--276} Cada elemento tem $12+3+3=18$ graus de liberdade.

\textbf{Quadratura.} Regra de 3 pontos interiores, $(\xi,\eta)\in\{(\tfrac16,\tfrac16),(\tfrac23,\tfrac16),(\tfrac16,\tfrac23)\}$,
peso $\tfrac16$, igual à das notas \nota{§14} e à do artigo \art{Seção~5.1}. \loc{UnsatSlope.h}{194} Ela integra
exatamente polinômios de grau 2, logo $\mathbf K$, $\mathbf Q$, $\mathbf C$ consistente, $\mathbf H$ e o peso com
coeficientes constantes; com coeficientes não lineares é aproximada.

\textbf{Condensação do armazenamento.} As notas propõem $\mathbf C^{L}_{\pi\beta}=\operatorname{diag}(\int C_{\pi\beta}N_i\,\dd\Omega)$
\nota{§14}. O código reconstrói as pressões nodais do P1 a partir do valor e do gradiente no ponto, avalia o termo de
armazenamento com as pressões de cada vértice e o multiplica por $\int N_i\,\dd\Omega$. Só os termos de saturação e de
compressibilidade do ar são condensados; o acoplamento e a compressibilidade da água ficam nos pontos de integração.
Para $C_{\pi\beta}$ constante as duas formas coincidem. \loc{TPZMatUnsaturatedUPwPa.cpp}{132--150, 235--240, 335--350}
\disc{D13}

%======================================================================
\section{Matrizes elementares e forma matricial}\label{sec:matriz}
%======================================================================

\nota{§13.4}, com a correspondência no código:
\begin{center}\small
\renewcommand{\arraystretch}{1.35}
\begin{tabular}{@{}p{0.17\linewidth}p{0.36\linewidth}p{0.41\linewidth}@{}}
\toprule
bloco & notas & código\\
\midrule
rigidez & $\mathbf K=\int\mathbf B_u^{\mathsf T}\mathbf D_T\mathbf B_u$ & $\int\mathbf B_u^{\mathsf T}\mathbb D\,\mathbf B_u$, com $\mathbb D$ da Seção~\ref{sec:mctan} (268--289)\\
$\mathbf u\leftarrow p_\pi$ & $\mathbf Q_{u\pi}=\int\mathbf B_u^{\mathsf T}\mathbf m\,\alpha\chi_\pi\mathbf N_p$ & $\lambda_p\,\mathbf Q_{u\pi}$ (290--294)\\
$p_\pi\leftarrow\mathbf u$ & $\mathbf Q_{\pi u}=\int\mathbf N_p^{\mathsf T}B_\pi\mathbf m^{\mathsf T}\mathbf B_u$ & igual, com $B_\pi=\alpha S_\pi$ (354--357)\\
armazenamento & $\mathbf C_{\pi\beta}=\int\mathbf N_p^{\mathsf T}C_{\pi\beta}\mathbf N_p$ & forma secante condensada (Seção~\ref{sec:temporal})\\
condução & $\mathbf H_\pi=\int\mathbf B_p^{\mathsf T}\lambda_\pi\mathbf B_p$ & $\int\mathbf B_p^{\mathsf T}\lambda_\pi\mathbf K\mathbf B_p$, $\lambda_\pi$ da Eq.~\eqref{eq:darcycod}\\
forças & $\mathbf f_u=\int\mathbf N_u^{\mathsf T}\rho_b\mathbf g+\int_{\Gamma_t}\mathbf N_u^{\mathsf T}\bar{\mathbf t}$ & $\lambda_g\int\mathbf N_u^{\mathsf T}\gamma\mathbf g+\mathbf f_t$\\
fluxos & $\mathbf f_\pi=\int\mathbf B_p^{\mathsf T}\lambda_\pi\rho_\pi\mathbf g-\int_{\Gamma_q}\mathbf N_p^{\mathsf T}\bar q_\pi$ & termo gravitacional dentro do resíduo; chuva nodal\\
densidade & $\mathbf g_\pi=\int\mathbf N_p^{\mathsf T}c_\pi\,\mathbf q_\pi\cdot\nabla p_\pi$ & só o ar\\
\bottomrule
\end{tabular}
\end{center}
Números entre parênteses: linhas de \path{TPZMatUnsaturatedUPwPa.cpp}.

\textbf{Sistema semidiscreto} \nota{§13.5}, com $\mathbf X=(\mathbf U,\mathbf P_w,\mathbf P_a)$:
\begin{equation}
\begin{bmatrix}\mathbf K&-\mathbf Q_{uw}&-\mathbf Q_{ua}\\
\mathbf Q_{wu}&\mathbf C_{ww}&\mathbf C_{wa}\\
\mathbf Q_{au}&\mathbf C_{aw}&\mathbf C_{aa}\end{bmatrix}
\begin{Bmatrix}\dot{\mathbf U}\\ \dot{\mathbf P}_w\\ \dot{\mathbf P}_a\end{Bmatrix}
+\begin{bmatrix}\mathbf 0&\mathbf 0&\mathbf 0\\ \mathbf 0&\mathbf H_w&\mathbf 0\\ \mathbf 0&\mathbf 0&\mathbf H_a\end{bmatrix}
\begin{Bmatrix}\mathbf U\\ \mathbf P_w\\ \mathbf P_a\end{Bmatrix}
=\begin{Bmatrix}\dot{\mathbf f}_u\\ \mathbf f_w-\mathbf g_w\\ \mathbf f_a-\mathbf g_a\end{Bmatrix}.
\end{equation}
O código não usa esta forma linearizada. Resolve o resíduo não linear em $t_{n+1}$, como na §13.7 das notas.

%======================================================================
\section{Discretização temporal e resíduos discretos}\label{sec:temporal}
%======================================================================

As notas usam o método $\theta$ e recomendam $\theta=1$ \nota{§13.6}; o resíduo em $t_{n+1}$ tem o armazenamento na
forma $(B_\pi\Delta\varepsilon_v+C_{\pi w}\Delta p_w+C_{\pi a}\Delta p_a)/\Delta t$ \nota{§13.7}. O código usa Euler
implícito e multiplica as equações de massa por $\Delta t$. Com $X=\alpha\Delta\varepsilon_v+\eta\Delta p_B$,
$\Delta(\cdot)=(\cdot)^{n+1}-(\cdot)^n$, todos os coeficientes em $n+1$ e $\langle\cdot\rangle$ o termo condensado
(Seção~\ref{sec:espacial}) \loc{TPZMatUnsaturatedUPwPa.cpp}{299--352}:
\begin{align}
\mathbf R_w&=\mathsf A_e\Big[\int\mathbf N_p^{\mathsf T}\Big(S_wX+\frac{nS_w}{K_w}\Delta p_w\Big)
+\big\langle n\,\Delta S_w\big\rangle
+\Delta t\int\mathbf B_p^{\mathsf T}\lambda_w\mathbf K(\nabla p_w-\gamma_w\mathbf g)\Big]-\Delta t\,\mathbf f_w,
\label{eq:Rw}\\
\mathbf R_a&=\mathsf A_e\Big[\int\mathbf N_p^{\mathsf T}\Big(S_aX-\Delta t\,c_a\lambda_a(\mathbf K\mathbf h_a)\cdot\nabla p_a\Big)
+\Big\langle -n\,\Delta S_w+\frac{nS_a}{\pabs}\Delta p_a\Big\rangle\notag\\
&\qquad\quad+\Delta t\int\mathbf B_p^{\mathsf T}\lambda_a\mathbf K\mathbf h_a\Big]-\Delta t\,\mathbf f_a,
\label{eq:Ra}
\end{align}
com $\mathbf h_a=\nabla p_a-\gamma_a\mathbf g$ e $\mathbf R_u$ da Eq.~\eqref{eq:Ru}. Na convenção do NeoPZ, o vetor
elementar é $-\mathbf R$ e a matriz é $\partial\mathbf R/\partial\mathbf X$. Com pressões congeladas (fator de segurança),
só $\mathbf R_u$ é montado. \loc{TPZMatUnsaturatedUPwPa.cpp}{297}

\textbf{Armazenamento da água.} Com $\eta=0$, as notas dão $(nS_w/K_w)\Delta p_w+n\Gamma^{n+1}(\Delta p_w-\Delta p_a)$, a
forma em taxa. O código usa $(nS_w/K_w)\Delta p_w+n(S_w^{n+1}-S_w^n)$, a forma secante. As duas coincidem quando
$\Delta t\to0$. A forma secante é a forma mista de Celia et al.\ (1990), referência [16] de Callari e Abati (2009,
p.~17). Callari e Abati (2009, p.~5) registram que a discretização da forma em taxa ``can lead to unacceptable errors in
fluid mass balance, especially in the case of strongly non-linear response, typically characterizing infiltration
problems''. \disc{D8}

\textbf{Armazenamento do ar.} A massa de ar por volume é $nS_a\rho_a$, com $\rho_a\propto\pabs$. Sua variação no passo,
dividida por $\rho_a^{n+1}$, é $n\Delta S_a+nS_a^{n}\Delta p_a/\pabs{}^{n+1}$. O código usa $S_a^{n+1}$ no segundo termo; a
diferença é $n\,\Delta S_a\,\Delta p_a/\pabs{}^{n+1}$, de segunda ordem no passo. No balanço global do modelo acoplado
com ar ativo e 36~mm/h (\path{Projects/SlopeRainfall/scripts/balanco_massa.py}, rodada
\path{resultados/artigo_rev/B_iv_f50_a45_I36_balanco}), o erro de água foi $\le0{,}05\%$ do volume infiltrado. O de
ar foi 0,25\% (24~h) e 0,34\% (48~h) do volume inicial de ar, ou 0,9--1,1\% do ar expulso. \disc{D9}

%======================================================================
\section{Jacobiano por blocos}\label{sec:jac}
%======================================================================

As notas montam o Jacobiano com os blocos $\mathbf G$ (peso dependente do estado) e $\tilde{\mathbf H}$ (variação da
mobilidade), e dispensam as derivadas de $B_\pi$ e $C_{\pi\beta}$ multiplicadas pelos incrementos (Newton modificado)
\nota{§13.7}. O código monta a derivada completa dos resíduos \eqref{eq:Ru}, \eqref{eq:Rw} e \eqref{eq:Ra}
\loc{TPZMatUnsaturatedUPwPa.cpp}{268--296, 354--378}. Com $\lambda'_\pi=\partial\lambda_\pi/\partial p_w=-\partial\lambda_\pi/\partial p_a$,
$\mathbf h_w=\nabla p_w-\gamma_w\mathbf g$ e $\mathrm{adv}=(\mathbf K\mathbf h_a)\cdot\nabla p_a$:
\begin{align*}
\mathbf J_{uu}&=\int\mathbf B_u^{\mathsf T}\mathbb D\,\mathbf B_u,\qquad
\mathbf J_{\pi u}=\int\mathbf N_p^{\mathsf T}S_\pi\,\alpha\,\mathbf m^{\mathsf T}\mathbf B_u,\\
\mathbf J_{u\pi}&=-\int\Big[\lambda_p\alpha\,\mathbf B_u^{\mathsf T}\mathbf m\,\chi_\pi+\lambda_g\mathbf N_u^{\mathsf T}\mathbf g\,\frac{\partial\gamma}{\partial p_\pi}\Big]\mathbf N_p,\\
\mathbf J_{ww}&=\int\mathbf N_p^{\mathsf T}\Big[\Gamma X+S_w\eta\chi_w+\frac{n}{K_w}(\Gamma\Delta p_w+S_w)\Big]\mathbf N_p
+\Delta t\int\mathbf B_p^{\mathsf T}\lambda_w\mathbf K\mathbf B_p\\
&\quad+\Delta t\int\mathbf B_p^{\mathsf T}(\mathbf K\mathbf h_w)\lambda'_w\mathbf N_p+\langle n\Gamma\rangle,\\
\mathbf J_{wa}&=\int\mathbf N_p^{\mathsf T}\Big[-\Gamma X+S_w\eta\chi_a-\frac{n}{K_w}\Gamma\Delta p_w\Big]\mathbf N_p
-\Delta t\int\mathbf B_p^{\mathsf T}(\mathbf K\mathbf h_w)\lambda'_w\mathbf N_p-\langle n\Gamma\rangle,\\
\mathbf J_{aw}&=\int\mathbf N_p^{\mathsf T}\Big[-\Gamma X+S_a\eta\chi_w-\Delta t\,c_a\lambda'_a\,\mathrm{adv}\Big]\mathbf N_p
+\Delta t\int\mathbf B_p^{\mathsf T}(\mathbf K\mathbf h_a)\lambda'_a\mathbf N_p
-\Big\langle n\Gamma\Big(1+\frac{\Delta p_a}{\pabs}\Big)\Big\rangle,\\
\mathbf J_{aa}&=\int\mathbf N_p^{\mathsf T}\Big[\Gamma X+S_a\eta\chi_a+\Delta t\,c_a^2\lambda_a\,\mathrm{adv}+\Delta t\,c_a\lambda'_a\,\mathrm{adv}\Big]\mathbf N_p\\
&\quad-\Delta t\int\mathbf N_p^{\mathsf T}c_a\lambda_a\big[\mathbf K(\mathbf h_a+\nabla p_a)\big]^{\mathsf T}\mathbf B_p\\
&\quad+\Delta t\int\mathbf B_p^{\mathsf T}\lambda_a\mathbf K\mathbf B_p
-\Delta t\int\mathbf B_p^{\mathsf T}(\mathbf K\mathbf h_a)\lambda'_a\mathbf N_p
+\Big\langle n\Gamma\Big(1+\frac{\Delta p_a}{\pabs}\Big)+\frac{nS_a}{\pabs}-\frac{nS_a\,\Delta p_a}{(\pabs)^2}\Big\rangle .
\end{align*}
O Jacobiano não é simétrico, como nas notas. Em relação às notas, o código acrescenta as derivadas dos coeficientes
vezes os incrementos (Newton completo) e não tem $\mathbf G_{uu}$, porque $n$ é constante \disc{D11}. O README do
programa do talude afirma convergência quadrática; a conferência por teste de Taylor do elemento acoplado é o teste~0 da
Fase~2.

%======================================================================
\section{Solução não linear}\label{sec:newton}
%======================================================================

As notas sugerem começar com Picard e passar a Newton \nota{§13.7}. O código usa Newton com o Jacobiano completo desde a
primeira iteração \loc{UnsatSlope.h}{331--410}:
\begin{itemize}
\item as condições de Dirichlet eliminam equações por filtro; o sistema é esparso e resolvido por LU
\loc{UnsatSlope.h}{244--251};
\item convergência quando $\|\mathbf R_u\|/\|\mathbf f_{\mathrm{peso}}\|<10^{-7}$ e $\max|\delta p|<10^{-6}$~kPa, depois de
pelo menos uma iteração; no máximo 40 iterações, ou 25 com pressões congeladas \loc{UnsatSlope.h}{381, 547--548};
\item busca linear (até quatro meias-passadas) sobre a norma do resíduo das pressões, enquanto a correção de pressão
for maior que a tolerância \loc{UnsatSlope.h}{396--406};
\item se o passo não converge, o passo de tempo é dividido ao meio \loc{main.cpp}{613--619}.
\end{itemize}
O resíduo das equações de massa não entra no critério; sua convergência é inferida da correção de pressão
\disc{D19}.

%======================================================================
\section{Fator de segurança: SRM e GIM}\label{sec:fs}
%======================================================================

Nem as notas nem o artigo tratam do fator de segurança. O código \loc{main.cpp}{443--516}:
\begin{itemize}
\item \textbf{Pressões congeladas} no instante $t$: só $\mathbf R_u$ é resolvido. O estado de tensões do transiente é
guardado e restaurado ao fim.
\item \textbf{Redução de resistência (SRM).} A partir do estado sem tensões, o peso é aplicado com resistência
aumentada, $F_0=0{,}6$, em quatro incrementos com as pressões \loc{main.cpp}{72, 421--431}. Depois $F$ cresce por
uma lista de valores e, na primeira não convergência, por bisseção até a tolerância de 0,01. O FS é o último $F$ que
convergiu. A não convergência em 25 iterações é o critério de ruptura. \loc{main.cpp}{167--192, 466--483}
\item \textbf{Aumento de gravidade (GIM).} A partir do estado sem tensões e com a resistência real, $\lambda$ multiplica
o peso e a parte positiva de $p_B$; a sucção fica com seu valor. Incrementos de 0,5 até 4, depois 4,5; 5; 6; 7; 8; 10;
bisseção até 0,01. \loc{main.cpp}{485--508}
\end{itemize}
Duas consequências. O FS não usa o estado de tensões do transiente acoplado, só as pressões dele \disc{D18}. E o fator
de gravidade não é o FS de redução de resistência: segundo o README de \path{Projects/SlopeMohrCoulomb}, multiplicar
$\gamma$ por $\lambda$ equivale a dividir só $c$ por $\lambda$, logo $FS_{\mathrm{GIM}}\ge FS_{\mathrm{SRM}}$ quando
$\phi>0$.

\textbf{Resultados já obtidos em solo seco} (README de \path{Projects/SlopeMohrCoulomb}, ciclo de refinamento 5,
$\nu=0{,}49$): FS por gravidade 1,793 e por SRM 1,203, contra 1,806 e 1,206 do método de Bishop simplificado.

%======================================================================
\section{Discrepâncias entre a formulação e o código}\label{sec:disc}
%======================================================================

``Notas'' e ``artigo'' são as fontes [N] e [A]. Na última coluna, \emph{manter} quer dizer que a simplificação é
aceitável e só precisa estar declarada; \emph{decidir} e \emph{corrigir} entram na Seção~\ref{sec:decisoes} e no plano.

{\small
\renewcommand{\arraystretch}{1.25}\setlength{\tabcolsep}{3pt}
\begin{longtable}{@{}p{0.05\linewidth}p{0.15\linewidth}p{0.21\linewidth}p{0.24\linewidth}p{0.29\linewidth}@{}}
\toprule
& tema & fontes & NeoPZ & efeito e proposta\\
\midrule
\endfirsthead
\toprule
& tema & fontes & NeoPZ & efeito e proposta\\
\midrule
\endhead
D1 & parâmetro de Bishop & $\chi=S_w$ \nota{§1, §3} & padrão $\chi_0=1$ ($\phi^b=\phi'$); opções $S_w$ e $S_e$ &
o equilíbrio usa $\chi_w=1$ e a massa usa $B_w=S_w$, par que as notas não derivam; a sucção dá mais resistência (FS
inicial do talude de referência 2,119 com $\chi=1$ e 1,895 com $\chi=S_e$, manuscrito,
\path{secao_generalizacao.tex}). \textbf{Decidir.}\\
D2 & Biot & $\alpha(n)=1-K/K_s$ & $\alpha=1$, $K_s\to\infty$, $\eta=0$ & grãos incompressíveis. \textbf{Manter.}\\
D3 & porosidade & $n(\mathbf z)$ linearizada \nota{§9} & $n=n_0$ nos coeficientes & erro da ordem de $\varepsilon_v$ em
$n\Delta S$ e $nS/K_w$. \textbf{Manter.}\\
D4 & permeabilidade intrínseca & Kozeny--Carman \nota{§5} & $\mathbf K$ constante (isotrópico ou anisotrópico) & sem
acoplamento da permeabilidade com a deformação. \textbf{Manter.}\\
D5 & peso & $\rho_b$ da mistura \nota{§4} & $\gamma=20$~kN/m$^3$ constante; mistura sem o ar como opção & no padrão o
peso não cresce com a infiltração. \textbf{Decidir} por teste (usar a mistura nos testes acoplados).\\
D6 & massa específica do ar & $\rho_a(p_a)$ & $\gamma_a$ constante no termo gravitacional & $\gamma_a\approx0{,}012$~kN/m$^3$:
efeito desprezível. \textbf{Manter.}\\
D7 & leis hidráulicas & van Genuchten e Mualem (água e ar) & padrão Fredlund--Xing, $S_e^4$ e Brooks--Corey;
van Genuchten e Mualem iguais às notas como opção & escolha de parâmetros, não de formulação. Usar a lei de cada
referência nos testes.\\
D8 & armazenamento da água & forma em taxa $C_{\pi\beta}\Delta p_\beta$ \nota{§13.7} & forma secante
$n(S_w^{n+1}-S_w^n)$ & conserva a massa de água (erro medido $\le0{,}05\%$). \textbf{Manter} e atualizar as notas.\\
D9 & armazenamento do ar & $C_{aa}\Delta p_a$ em taxa & $n\Delta S_a+nS_a^{n+1}\Delta p_a/\pabs$ & erro de
segunda ordem; balanço de ar 0,25--0,34\%. \textbf{Corrigir} para $S_a^{n}$ antes do teste~3.\\
D10 & termo $c_w\mathbf q_w\cdot\nabla p_w$ & opcional \nota{§13.4} & ausente & as notas o dispensam. Nada a fazer.\\
D11 & Jacobiano & Newton modificado \nota{§13.7} & Newton completo, sem $\mathbf G_{uu}$ & deve dar convergência
quadrática. \textbf{Verificar} (teste~0).\\
D12 & interpolação & mesma ordem, Q4 ou T3 \nota{§13.3, §14} & Taylor--Hood T6/T3 & segue a recomendação das próprias
notas \nota{§13.9}. Atualizar a §14 das notas para o T6.\\
D13 & condensação & $\operatorname{diag}(\int C N_i)$ & valor nodal vezes $\int N_i$, só saturação e compressibilidade
do ar & igual para coeficiente constante. \textbf{Manter.}\\
D14 & Mohr--Coulomb não associado & artigo associado & PV2 só associado; $\psi<\phi$ só no \code{SlopeRainfallPV}, com a
classe antiga & os resultados com $\psi<\phi$ do manuscrito não vêm do modelo do artigo. \textbf{Decidir.}\\
D15 & redução de resistência & --- & $c/F$, $\tan\phi/F$, $\tan\psi/F$ & documentada na Seção~\ref{sec:mc}.\\
D16 & estado inicial & $\sig_0$ e $p_{B0}$ dados & geostático pelo modelo; sucção limitada a 75~kPa & o campo de
pressões não é estacionário em $t=0$. Medir a deriva sem chuva no teste~4.\\
D17 & contorno do ar & --- & $p_a=0$ na superfície mesmo alagada & o aprisionamento do ar depende de $k_{ra}^{\min}$
na camada saturada. Testar a sensibilidade a $k_{ra}^{\min}$.\\
D18 & fator de segurança & --- & SRM a partir de $F_0$ e GIM sem escalar a sucção, com pressões congeladas & o FS não
usa o estado de tensões do transiente. Testar a independência de $F_0$ (teste~1).\\
D19 & critério de convergência & --- & $\|\mathbf R_u\|$ relativo e $\max|\delta p|$ & o resíduo de massa não é testado
diretamente. Acrescentar a norma do resíduo de massa ao relatório de convergência.\\
D20 & sinal do fluxo prescrito & $\bar q>0$ saindo & fluxo positivo entrando & só convenção.\\
\bottomrule
\end{longtable}
}

%======================================================================
\section{Decisões antes da Fase 2}\label{sec:decisoes}
%======================================================================

\begin{enumerate}
\item \textbf{$\chi$ dos testes (D1).} Proposta: $\chi=S_w$, como nas notas, em todos os testes de formulação. O
$\chi=1$ fica só para a comparação com Rahardjo et al.\ (2007).
\item \textbf{Armazenamento do ar (D9).} Proposta: trocar $S_a^{n+1}$ por $S_a^{n}$ no termo de compressibilidade, com o
Jacobiano correspondente, e medir antes e depois com o teste de Boyle (teste~3a).
\item \textbf{Regra não associada (D14).} O PV2 do artigo é associado. Opções: (a) restringir a análise a $\psi=\phi$
com o PV2; (b) estender a projeção do artigo ao caso não associado, que exige outra derivação, porque a projeção
deixa de ser a de ponto mais próximo na norma de energia; (c) manter a classe antiga só como referência.
Proposta: (a), já que o modelo escolhido é o PV2.
\end{enumerate}

%======================================================================
\section{Plano de verificação da Fase 2}\label{sec:plano}
%======================================================================

Um teste por vez, com critério de aceitação fixado antes de rodar, e um commit por teste com tabela ou figura.

\textbf{Teste 0. Jacobiano do elemento acoplado.} Perturbar $\mathbf X$ numa direção aleatória e medir
$E(\alpha)=\|\mathbf R(\mathbf X+\alpha\,\delta\mathbf X)-\mathbf R(\mathbf X)-\alpha\,\mathbf J\,\delta\mathbf X\|$, como no
teste de Taylor do artigo \art{Eqs.~63--65}. Estados: saturado, não saturado, plástico; $\chi=S_w$ e $\chi=1$; com e sem
condensação. \emph{Aceitação:} ordem $p\in[1{,}9;\,2{,}1]$ em todos. \emph{Situação:} a implementar.

\textbf{Teste 1. Solo seco: GIM e SRM, sem água e sem ar.}
\begin{itemize}
\item 1a, só mecânica (\path{Projects/SlopeMohrCoulomb}): feito; ciclo 5 com 1,793 e 1,203 contra Bishop 1,806 e 1,206.
Falta verificar a independência do FS em relação a $F_0$.
\item 1b, o mesmo talude no programa acoplado com $p_w=p_a=0$ congeladas: deve reproduzir 1a na mesma malha.
\item \emph{Aceitação:} diferença entre 1a e 1b $\le0{,}01$; variação do FS com $F_0\in\{0{,}5;\,0{,}6;\,0{,}8\}$
$\le0{,}01$; convergência com o refinamento para o valor de Bishop dentro de 1\%.
\end{itemize}

\textbf{Teste 2. Água com ar passivo e esqueleto rígido (equação de Richards).}
\begin{itemize}
\item Infiltração de Philip conforme Callari e Abati (2009, Seção~4.2, pp.~8--9): leis de van Genuchten (Eqs.~59--60),
parâmetros da Tabela~2 (p.~8), malhas $2\times6\times60$ e $2\times2\times20$, $\Delta t=1$~s nos 1000 primeiros passos e
10~s depois, perfil de pressão comparado em $t=21\,600$~s (Figs.~6--7).
\item No código: deslocamentos fixos, sem acoplamento e com $p_a=0$ em todo o domínio.
\item \emph{Aceitação:} frente de saturação dentro de um elemento da solução semianalítica; erro de balanço de água
$\le0{,}1\%$; resultado estável com $\Delta t/2$ e com a malha refinada.
\end{itemize}

\textbf{Teste 3. Água e ar com esqueleto rígido.}
\begin{itemize}
\item 3a, teste de Boyle: coluna rígida com o ar fechado em todo o contorno e água injetada. A massa de ar se conserva
e, com gás ideal isotérmico, $\pabs\,V_a$ deve ficar constante, com $V_a=\int nS_a\,\dd\Omega$. \emph{Aceitação:} variação
relativa de $\pabs V_a$ $\le10^{-4}$ depois da correção D9. Mede diretamente o efeito de D9.
\item 3b, limite passivo: com o ar livre nas duas extremidades, a solução com ar ativo deve tender à de ar passivo
quando $\mu_a\to0$. \emph{Aceitação:} a diferença em $p_w$ entre as duas soluções cai quando $\mu_w/\mu_a$ cresce e,
com $\mu_w/\mu_a$ cem vezes o valor físico, fica abaixo de 1\% da variação de $p_w$ no ensaio.
\item 3c, balanço global de água e ar em uma infiltração com alagamento. \emph{Aceitação:} erro de ar $\le0{,}1\%$ do ar
inicial depois de D9.
\end{itemize}

\textbf{Teste 4. Três campos acoplados.}
\begin{itemize}
\item 4a, Terzaghi no limite saturado ($S_w=1$): comparar com a solução em série já usada no projeto
\path{Projects/TerzaghiConsolidation} (outro material). \emph{Aceitação:} erro máximo de $p_w/p_0\le1\%$ na malha e no passo
refinados.
\item 4b, Mandel, saturado: solução analítica do Apêndice~A.1 de Haghighat et al.\ (2021, p.~25), com as raízes de
$\tan\alpha_i=\alpha_i(1-\nu)/(\nu_u-\nu)$. Verifica o efeito Mandel--Cryer com P2/P1. \emph{Aceitação:} erro de $p$ no
centro $\le2\%$.
\item 4c, Liakopoulos com ar passivo contra os dados experimentais: Callari e Abati (2009, Tabela~1, Eq.~58, Fig.~3:
vazão na base e pressões a 600, 1200 e 1800~s). \emph{Aceitação:} desvio da mesma ordem que o da solução de Callari e
Abati contra os dados.
\item 4d, Liakopoulos com ar ativo: configuração de Haghighat et al.\ (2021, Seção~5.3, p.~23), com gás à pressão
atmosférica no topo e na base. A referência deles é numérica, então a comparação é qualitativa.
\item 4e, deriva do estado inicial do talude sem chuva (D16).
\end{itemize}

%======================================================================
\section{Referências citadas}\label{sec:refs}
%======================================================================

\begin{itemize}\small
\item [N] Formulação poroelástica não saturada: esqueleto deformável, água e ar. Notas,
\path{docs/biot_nao_saturado/biot_nao_saturado_agua_ar3.tex}.
\item [A] Lira Cecílio, D., 2026. Analytical return mapping for Mohr--Coulomb plasticity in rotated Haigh--Westergaard
space with consistent tangent operator. Journal of Engineering Mathematics 157, 10. doi:10.1007/s10665-026-10519-1.
\item Callari, C., Abati, A., 2009. Finite element methods for unsaturated porous solids and their application to dam
engineering problems. Computers \& Structures 87, 485--501.
\item Celia, M.A., Bouloutas, E.T., Zarba, R.L., 1990. A general mass-conservative numerical solution for the
unsaturated flow equation. Water Resources Research 26(7), 1483--1496. Citada pela referência [16] de Callari e Abati
(2009, p.~17); o PDF não foi lido.
\item Haghighat, E., Amini, D., Juanes, R., 2021. Physics-informed neural network simulation of multiphase
poroelasticity using stress-split sequential training. arXiv:2110.03049v1.
\item Pedroso, D.M., 2015. A solution to transient seepage in unsaturated porous media. Computer Methods in Applied
Mechanics and Engineering 285, 791--816.
\item Rahardjo, H., Ong, T.H., Rezaur, R.B., Leong, E.C., 2007. Factors controlling instability of homogeneous soil
slopes under rainfall. Journal of Geotechnical and Geoenvironmental Engineering 133(12), 1532--1543.
\end{itemize}

A revisão bibliográfica com os 43 PDFs enviados é a próxima parte deste documento.

\appendix
%======================================================================
\section{Arquivos da verificação}\label{sec:arquivos}
%======================================================================

Em \path{docs/biot_nao_saturado/formulacao/verificacao/}:
\begin{itemize}
\item \path{verifica_pv2.py} e \path{verifica_pv2.txt}: expressões do PV2 contra as Tabelas~4--6 do artigo, contra
diferenças finitas e contra a projeção de ponto mais próximo por conjunto ativo. Uso:
\texttt{python3 verifica\_pv2.py Material/Plasticity/TPZYCMohrCoulombPV2.cpp}.
\item \path{taylor_check.txt}: saída de \code{SlopeMohrCoulomb check}, teste de Taylor da tangente.
\item \path{eq54.py} e \path{eq54.txt}: razão entre a Eq.~(54) do artigo e o quadrado da distância.
\end{itemize}

\end{document}
